\documentclass[11pt]{article}

\usepackage[T1]{fontenc}
\usepackage{amsmath,amssymb,amsthm,mathtools}
\usepackage[margin=1in]{geometry}
\usepackage{booktabs}
\usepackage{hyperref}
\hypersetup{hidelinks}

\newtheorem{theorem}{Theorem}[section]
\newtheorem{lemma}[theorem]{Lemma}
\newtheorem{corollary}[theorem]{Corollary}
\newtheorem{proposition}[theorem]{Proposition}
\theoremstyle{remark}
\newtheorem{remark}[theorem]{Remark}

\newcommand{\R}{\mathcal R}

\title{Transported Romberg Extrapolation for the Arcsine Schr\"oder Linearizer\\
\large Polynomial Composition Lifts and a Flat Infinite-Depth Limit}
\author{Wayne Baker\\Independent Researcher}
\date{October 1, 2026}

\begin{document}
\maketitle

\begin{abstract}
For an odd integer $m\ge3$, let
\[
S_m(z)=\sin(m\arcsin z),
\qquad A(z)=\arcsin z,
\]
so that $S_m$ is an odd polynomial and $A(S_m(z))=mA(z)$ as an analytic germ at
$0$.  Starting from $P^{(m)}_1(z)=z$, consider the finite-polynomial recursion
\[
P^{(m)}_{j+1}
=P^{(m)}_j+
\frac{mP^{(m)}_j-P^{(m)}_j\circ S_m}{m^{2j+1}-m}.
\]
After the change of variable $z=\sin u$, this recursion is exactly a
Richardson--Romberg polynomial extrapolation.  It extrapolates
$T_u(h)=\sin(hu)/h$ to $h=0$ from the geometric nodes
$h=1,m,\ldots,m^{j-1}$.  Consequently
\[
P^{(m)}_j(\sin u)=\sum_{i=0}^{j-1}w^{(m)}_{j,i}\frac{\sin(m^i u)}{m^i},
\]
with the standard Lagrange weights at zero in the variable $h^2$.  The familiar
jet-recovery and leading-defect formulas are therefore classical extrapolation
consequences; in particular
\[
\arcsin z-P^{(m)}_j(z)
=\frac{m^{j(j-1)}}{(2j+1)!}z^{2j+1}+O(z^{2j+3}).
\]
The new focus is the infinite-depth limit.  The weights converge to
\[
w^{(m)}_{\infty,i}
=\frac{(-1)^i q^{i(i+1)/2}}
{(q;q)_i(q;q)_\infty},\qquad q=m^{-2},
\]
which decay super-geometrically.  Hence the limiting trigonometric series
\[
F_m(u)=\sum_{i\ge0}w^{(m)}_{\infty,i}
\frac{\sin(m^i u)}{m^i}
\]
is $C^\infty$ on $\mathbb R$, and the finite tableaux converge to it together
with every derivative.  All moments required by polynomial extrapolation
survive the limit, so $u-F_m(u)$ is flat at $u=0$.  The series is Hadamard-lacunary.  Its Fourier coefficients decay
super-algebraically but not exponentially; combined with the
functional-differential identity
\[
F_m(mu)=m\bigl(F_m''(u)+F_m(u)\bigr),
\]
this shows that $F_m$ is nowhere real analytic.  The transported function
$G_m=F_m\circ\arcsin$ is $C^\infty$ on $(-1,1)$ and differs from $\arcsin$
by a nonzero flat function at the origin.  The cubic case $m=3$ gives the
scale-three specialization.
\end{abstract}

\section{Introduction}

Schr\"oder's functional equation
\[
\Psi(f(z))=\lambda\Psi(z)
\]
is a classical framework for local linearization near a fixed point, going
back to Schr\"oder's nineteenth-century work \cite{Schroeder1870}.  Richardson
and Gaunt's deferred approach to the limit \cite{RichardsonGaunt1927}, and the
geometric-grid formulation associated with Romberg \cite{Romberg1955}, provide
an equally classical mechanism for cancelling successive powers in an
asymptotic expansion; see also Joyce's survey \cite{Joyce1971} and the monograph
of Brezinski and Redivo-Zaglia \cite{BrezinskiRedivo1991}.  A constructive
computational treatment of Schr\"oder linearization was given by Iserles
\cite{Iserles1993}, who also related the analysis to convergence acceleration of
functional iteration.  For a systematic modern treatment of Richardson
extrapolation, including geometric and near-geometric extrapolation sequences and
Romberg-type constructions, see Sidi \cite{Sidi2003}.

The present paper concerns the intersection of these two classical structures.
For every odd integer $m\ge3$, define
\begin{equation}
S_m(z):=\sin(m\arcsin z).
\label{eq:Sm}
\end{equation}
The multiple-angle formula makes $S_m$ an odd polynomial of degree $m$, and
near the origin
\begin{equation}
A(S_m(z))=mA(z),
\qquad A(z):=\arcsin z,
\qquad A(0)=0,\quad A'(0)=1.
\label{eq:schroeder}
\end{equation}
For $m=3$ this is the cubic sine-multiple map
\[
S_3(z)=3z-4z^3.
\]

Starting from $P^{(m)}_1(z)=z$, consider the recursion
\begin{equation}
P^{(m)}_{j+1}(z)
=
P^{(m)}_j(z)+
\frac{mP^{(m)}_j(z)-P^{(m)}_j(S_m(z))}
     {m^{2j+1}-m}.
\label{eq:lift}
\end{equation}
After writing $z=\sin u$, the recursion is precisely the Romberg/Neville
tableau obtained by extrapolating
\[
T_u(h):=\frac{\sin(hu)}{h}
\]
to $h=0$ through the nodes $h=m^i$.  Thus the finite-stage jet recovery and its
leading error constant are not a new extrapolation principle: they are the
standard polynomial extrapolation mechanism transported through the nonlinear
coordinate $z=\sin u$.

This identification leaves a more distinctive question.  The geometric nodes
move away from zero, rather than toward it, and the resulting Lagrange weights
have a super-geometric limiting profile.  Does the transported finite-polynomial
tableau converge as its depth tends to infinity, and if it does, what is the
limit?  We prove that the tableaux converge in $C^\infty$ in the linearizing
variable to a periodic smooth function whose Taylor series at the origin is
exactly $u$, but which is nowhere real analytic.  The limit is a
Hadamard-lacunary trigonometric series, placing it in a classical family studied
in harmonic analysis; see Zygmund~\cite{Zygmund2002}, Kahane~\cite{Kahane1964},
and Katznelson~\cite{Katznelson2004}.  The distinctive feature here is that the
full Taylor jet is prescribed exactly by the extrapolation moments.  After
transport back to $z$, the limiting polynomial sequence differs from $\arcsin z$
by a nonzero flat function.

Accordingly, the contribution claimed here is not Richardson/Romberg
extrapolation itself, nor constructive approximation of Schr\"oder linearizers,
nor the multiple-angle identity.  It is the exact transport of the classical
geometric extrapolation tableau through the arcsine Schr\"oder coordinate into
a hierarchy of forward polynomial compositions, together with the explicit
analysis of the resulting infinite-depth limit.  In particular, the limiting
object is smooth and has the full Taylor jet of the linearizer at the fixed
point, yet differs from it by a nonzero flat function.

\section{The transported Romberg tableau}
\label{sec:romberg}

Put
\[
x_i:=m^{2i},\qquad i=0,1,\ldots,
\]
and define the Lagrange weights for evaluation at $x=0$ on the first $j$ nodes
by
\begin{equation}
w^{(m)}_{j,i}
:=
\prod_{\substack{0\le k<j\\k\ne i}}
\frac{x_k}{x_k-x_i},
\qquad 0\le i<j.
\label{eq:finiteweights}
\end{equation}
For fixed $u$, set
\begin{equation}
\mathcal T_u(x)
:=
\begin{cases}
\dfrac{\sin(u\sqrt{x})}{\sqrt{x}},&x\ne0,\\[1ex]
u,&x=0.
\end{cases}
\label{eq:Tux}
\end{equation}
Its power series
\[
\mathcal T_u(x)
=
\sum_{r=0}^{\infty}
\frac{(-1)^r u^{2r+1}}{(2r+1)!}x^r
\]
shows that it is entire in $x$.

\begin{theorem}[Transported Romberg representation]
\label{thm:representation}
For every odd $m\ge3$ and every $j\ge1$,
\begin{equation}
P^{(m)}_j(\sin u)
=
\sum_{i=0}^{j-1}w^{(m)}_{j,i}
\frac{\sin(m^iu)}{m^i}
=
\sum_{i=0}^{j-1}w^{(m)}_{j,i}\,\mathcal T_u(x_i).
\label{eq:representation}
\end{equation}
Equivalently, $P^{(m)}_j(\sin u)$ is the value at $x=0$ of the degree-$(j-1)$
polynomial interpolating $\mathcal T_u(x)$ at
$x=1,m^2,\ldots,m^{2(j-1)}$.
In the $z$ variable,
\begin{equation}
P^{(m)}_j(z)
=
\sum_{i=0}^{j-1}w^{(m)}_{j,i}\,m^{-i}
S_m^{\circ i}(z).
\label{eq:iterateform}
\end{equation}
\end{theorem}

\begin{proof}
The interpolation formula gives the second expression in
\eqref{eq:representation}.  Since $x_i=m^{2i}$,
$\mathcal T_u(x_i)=\sin(m^iu)/m^i$.

It remains to identify the interpolation tableau with \eqref{eq:lift}.
Writing
$\widehat P^{(m)}_j(u):=P^{(m)}_j(\sin u)$ and using
$S_m(\sin u)=\sin(mu)$, the lift becomes
\begin{equation}
\widehat P^{(m)}_{j+1}(u)
=
\frac{m^{2j+1}\widehat P^{(m)}_j(u)
      -\widehat P^{(m)}_j(mu)}{m^{2j+1}-m}
=
\frac{m^{2j}\widehat P^{(m)}_j(u)
      -m^{-1}\widehat P^{(m)}_j(mu)}{m^{2j}-1}.
\label{eq:nevillerec}
\end{equation}
To identify this recursion with the interpolation tableau, note first the
scaling identity
\begin{equation}
\mathcal T_u(m^2x)=m^{-1}\mathcal T_{mu}(x).
\label{eq:scalingTu}
\end{equation}
Hence the degree-$(j-1)$ interpolant of $\mathcal T_u$ on the shifted node
set $x_1,\ldots,x_j$, evaluated at $0$, is
$m^{-1}\widehat P^{(m)}_j(mu)$ whenever
$\widehat P^{(m)}_j(u)$ is the corresponding interpolated value on
$x_0,\ldots,x_{j-1}$.  Neville's two-block formula at $x=0$, with
$x_0=1$ and $x_j=m^{2j}$, therefore gives
\[
\frac{x_j\widehat P^{(m)}_j(u)
      -m^{-1}\widehat P^{(m)}_j(mu)}
     {x_j-x_0}
=
\frac{m^{2j}\widehat P^{(m)}_j(u)
      -m^{-1}\widehat P^{(m)}_j(mu)}
     {m^{2j}-1},
\]
which is exactly \eqref{eq:nevillerec}.  Since
$\widehat P^{(m)}_1(u)=\sin u=\mathcal T_u(1)$ for all real $u$, induction
proves the Lagrange formula \eqref{eq:representation} for all real $u$.

Finally, the polynomial identity
$S_m^{\circ i}(\sin u)=\sin(m^iu)$ holds near $u=0$ and therefore identically
as a polynomial in $z=\sin u$, which gives \eqref{eq:iterateform}.
\end{proof}

\begin{remark}
The theorem makes the classical content explicit.  The finite polynomial lift
is not merely analogous to Richardson extrapolation: after the coordinate
change $z=\sin u$, it is exactly the same polynomial extrapolation tableau.
The nonlinear composition structure appears when that tableau is transported
back to $z$.
\end{remark}

\section{Finite-depth consequences}
\label{sec:finite}

The classical exactness of polynomial interpolation immediately recovers the
finite-order results.

\begin{proposition}[Moment identities]
\label{prop:finite-moments}
For $0\le r\le j-1$,
\begin{equation}
\sum_{i=0}^{j-1}w^{(m)}_{j,i}x_i^r
=
\begin{cases}
1,&r=0,\\
0,&1\le r\le j-1.
\end{cases}
\label{eq:finite-moments}
\end{equation}
Moreover,
\begin{equation}
\sum_{i=0}^{j-1}w^{(m)}_{j,i}x_i^j
=(-1)^{j-1}\prod_{i=0}^{j-1}x_i
=(-1)^{j-1}m^{j(j-1)}.
\label{eq:leadingmoment}
\end{equation}
\end{proposition}

\begin{proof}
The identities in \eqref{eq:finite-moments} follow from exact Lagrange
interpolation of $1,x,\ldots,x^{j-1}$ at $x=0$.  For $x^j$, let $I_{j-1}$
denote its interpolating polynomial.  The difference
$x^j-I_{j-1}(x)$ is monic and vanishes at all nodes $x_i$, so
\[
x^j-I_{j-1}(x)=\prod_{i=0}^{j-1}(x-x_i).
\]
Setting $x=0$ gives \eqref{eq:leadingmoment}.
\end{proof}

\begin{theorem}[Jet recovery and exact leading defect]
\label{thm:finite-defect}
For every odd $m\ge3$ and every $j\ge1$,
\begin{equation}
P^{(m)}_j(z)
=\arcsin z+O(z^{2j+1}),
\label{eq:jet}
\end{equation}
and more precisely
\begin{equation}
\arcsin z-P^{(m)}_j(z)
=
\frac{m^{j(j-1)}}{(2j+1)!}z^{2j+1}
+O(z^{2j+3}).
\label{eq:defectlaw}
\end{equation}
Also
\begin{equation}
\deg P^{(m)}_j=m^{j-1}.
\label{eq:degree}
\end{equation}
\end{theorem}

\begin{proof}
Insert the power series of $\mathcal T_u$ into \eqref{eq:representation}.
The moment identities annihilate the terms $x^r$ for $1\le r\le j-1$, while
\eqref{eq:leadingmoment} gives
\[
P^{(m)}_j(\sin u)
=
u-
\frac{m^{j(j-1)}}{(2j+1)!}u^{2j+1}
+O(u^{2j+3}).
\]
Thus $u-P^{(m)}_j(\sin u)$ has the displayed leading coefficient.  Since
$u=\arcsin z=z+O(z^3)$, transport back to $z$ gives
\eqref{eq:jet}--\eqref{eq:defectlaw}.

For the degree, \eqref{eq:iterateform} has a nonzero coefficient on
$S_m^{\circ(j-1)}$, whose degree is $m^{j-1}$, while all earlier iterates have
smaller degree.
\end{proof}

For the cubic case $m=3$, the first residual coefficients are therefore
\[
\frac{3^2}{5!}=\frac3{40},\qquad
\frac{3^6}{7!}=\frac{81}{560},\qquad
\frac{3^{12}}{9!}=\frac{6561}{4480},
\]
for $P_2,P_3,P_4$, respectively.  The first explicit lifts are
\begin{align}
P_2(z)&=z+\frac16z^3,\label{eq:P2}\\
P_3(z)&=z+\frac16z^3+\frac3{40}z^5
-\frac1{10}z^7+\frac2{45}z^9.\label{eq:P3}
\end{align}
The degree-$27$ fourth lift begins
\begin{equation}
P_4(z)
=z+\frac16z^3+\frac3{40}z^5+\frac5{112}z^7
-\frac{1807}{1260}z^9+O(z^{11}),
\label{eq:P4}
\end{equation}
where the $O(z^{11})$ notation suppresses the remaining terms of the finite
polynomial.

\section{Intertwining with a geometric scale tower}
\label{sec:tower}

The same tableau can be read directly in the scale variable.  Suppose
\begin{equation}
N_{k+1}=mN_k,
\qquad
s_k=S_m(s_{k+1}),
\label{eq:tower}
\end{equation}
with the points lying in a local branch on which \eqref{eq:schroeder} applies.
Then
\begin{equation}
N_kA(s_k)=N_{k+1}A(s_{k+1}),
\label{eq:invariant}
\end{equation}
so $N_k\arcsin s_k$ is exactly invariant along the tower.

For $j\ge1$, define
\[
F_j(k):=N_kP^{(m)}_j(s_k)
\]
and the level-$j$ Richardson operator
\begin{equation}
(\R_jF)(k)
:=
\frac{m^{2j}F(k+1)-F(k)}{m^{2j}-1}.
\label{eq:Rj}
\end{equation}

\begin{theorem}[Scale/composition intertwining identity]
\label{thm:intertwining}
Whenever \eqref{eq:tower} holds,
\begin{equation}
(\R_jF_j)(k)=F_{j+1}(k+1).
\label{eq:intertwining}
\end{equation}
\end{theorem}

\begin{proof}
Using $N_{k+1}=mN_k$ and $s_k=S_m(s_{k+1})$,
\begin{align*}
(\R_jF_j)(k)
&=
\frac{m^{2j}N_{k+1}P^{(m)}_j(s_{k+1})
      -N_kP^{(m)}_j(s_k)}{m^{2j}-1}\\
&=N_k
\frac{m^{2j+1}P^{(m)}_j(s_{k+1})
      -P^{(m)}_j(S_m(s_{k+1}))}{m^{2j}-1}\\
&=mN_kP^{(m)}_{j+1}(s_{k+1})\\
&=F_{j+1}(k+1).
\end{align*}
\end{proof}

For $m=3$ and $j=1$, writing $G_k=N_ks_k$ gives
\[
\frac{9G_{k+1}-G_k}{8}
=N_{k+1}\left(s_{k+1}+\frac{s_{k+1}^3}{6}\right),
\]
the familiar first scale-three correction.  The terminology
``intertwining'' is used deliberately: the identity has an index shift and is
not a conjugacy in the strict dynamical-systems sense.

\section{Infinite-depth weights}
\label{sec:weights}

We now pass from the classical finite tableau to the infinite-depth problem.
Put
\begin{equation}
q:=m^{-2}\in(0,1),
\qquad
(q;q)_n:=\prod_{r=1}^{n}(1-q^r),
\qquad
(q;q)_\infty:=\prod_{r=1}^{\infty}(1-q^r).
\label{eq:qpoch}
\end{equation}
The infinite product is strictly positive.

\begin{proposition}[Closed finite and limiting weights]
\label{prop:weights}
The Lagrange weights \eqref{eq:finiteweights} satisfy
\begin{equation}
w^{(m)}_{j,i}
=
\frac{(-1)^i q^{i(i+1)/2}}
{(q;q)_i\,(q;q)_{j-1-i}},
\qquad 0\le i<j.
\label{eq:finiteqweights}
\end{equation}
For every fixed $i$,
\begin{equation}
w^{(m)}_{j,i}\longrightarrow
w^{(m)}_{\infty,i}
:=
\frac{(-1)^i q^{i(i+1)/2}}
{(q;q)_i\,(q;q)_\infty}.
\label{eq:limitweights}
\end{equation}
Moreover, uniformly in $j>i$,
\begin{equation}
|w^{(m)}_{j,i}|
\le
\frac{q^{i(i+1)/2}}{(q;q)_\infty^2}.
\label{eq:weightbound}
\end{equation}
\end{proposition}

\begin{proof}
Since $x_i=q^{-i}$,
\[
w^{(m)}_{j,i}
=\prod_{k\ne i}\frac1{1-q^{k-i}}.
\]
For $k<i$, put $r=i-k$ and use
$1-q^{-r}=-q^{-r}(1-q^r)$; for $k>i$, put $r=k-i$.
The two products give \eqref{eq:finiteqweights}.  Letting $j\to\infty$ yields
\eqref{eq:limitweights}.  Finally
$(q;q)_n\ge(q;q)_\infty>0$ gives \eqref{eq:weightbound}.
\end{proof}

The exponent $i(i+1)/2$ shows that the limiting weights decay faster than any
geometric sequence.  This compensates for the increasing frequencies $m^i$.

\section{The smooth flat limit}
\label{sec:limit}

Define
\begin{equation}
F_m(u)
:=
\sum_{i=0}^{\infty}w^{(m)}_{\infty,i}
\frac{\sin(m^iu)}{m^i}.
\label{eq:Fm}
\end{equation}

\begin{theorem}[$C^\infty$ convergence of the infinite tableau]
\label{thm:Cinf}
For each integer $r\ge0$, the series obtained by differentiating
\eqref{eq:Fm} $r$ times converges absolutely and uniformly on $\mathbb R$.
Consequently $F_m\in C^\infty(\mathbb R)$, and
\begin{equation}
P^{(m)}_j(\sin u)\longrightarrow F_m(u)
\label{eq:Cinfconv}
\end{equation}
with uniform convergence of every derivative on $\mathbb R$.
\end{theorem}

\begin{proof}
Extend $w^{(m)}_{j,i}$ by zero for $i\ge j$.  For fixed $i$,
$w^{(m)}_{j,i}\to w^{(m)}_{\infty,i}$.  For $r=0$,
\[
\left|w^{(m)}_{j,i}\frac{\sin(m^iu)}{m^i}\right|
\le C q^{i(i+1)/2}m^{-i}.
\]
For $r\ge1$, the $r$th derivative of
$\sin(m^iu)/m^i$ has absolute value at most $m^{i(r-1)}$, and hence
\[
\left|w^{(m)}_{j,i}
\frac{d^r}{du^r}\frac{\sin(m^iu)}{m^i}\right|
\le
C q^{i(i+1)/2}m^{i(r-1)}
=Cm^{-i^2+(r-2)i},
\]
where $C=(q;q)_\infty^{-2}$.  The majorant is summable for every fixed $r$.
The Weierstrass test and Tannery's theorem (equivalently, dominated convergence
for these coefficient arrays) give uniform convergence of the series and of
every differentiated series.
\end{proof}

The finite interpolation moments survive the limit.

\begin{proposition}[Infinite moment identities]
\label{prop:infinite-moments}
For every integer $r\ge0$,
\begin{equation}
\sum_{i=0}^{\infty}w^{(m)}_{\infty,i}m^{2ri}
=
\begin{cases}
1,&r=0,\\
0,&r\ge1.
\end{cases}
\label{eq:infinite-moments}
\end{equation}
\end{proposition}

\begin{proof}
Fix $r$.  For all $j\ge r+1$, Proposition~\ref{prop:finite-moments} gives the
corresponding finite identity.  The bound
\[
|w^{(m)}_{j,i}|m^{2ri}
\le Cq^{i(i+1)/2}q^{-ri}
\]
is summable in $i$.  Tannery's theorem therefore allows $j\to\infty$ and
gives \eqref{eq:infinite-moments}.
\end{proof}

\begin{remark}[Euler's $q$-series form of the moments]
The same identities follow directly from the classical Euler
$q$-binomial identity
\[
\sum_{i=0}^{\infty}
\frac{(-1)^i q^{i(i-1)/2}y^i}{(q;q)_i}
=(y;q)_\infty,
\qquad |q|<1,
\]
see, for example, Gasper and Rahman~\cite{GasperRahman2004}.  Indeed,
using $m^{2r}=q^{-r}$ and \eqref{eq:limitweights},
\[
\sum_{i=0}^{\infty}w^{(m)}_{\infty,i}m^{2ri}
=
\frac{(q^{1-r};q)_\infty}{(q;q)_\infty}.
\]
For $r=0$ this ratio is $1$, while for every $r\ge1$ the numerator contains
the factor $1-q^0$ and therefore vanishes.  This gives an independent
$q$-series check of Proposition~\ref{prop:infinite-moments}.
\end{remark}

\begin{theorem}[Flat and nowhere-analytic infinite-depth limit]
\label{thm:flat}
The function $F_m$ satisfies
\begin{equation}
F_m'(0)=1,
\qquad
F_m^{(r)}(0)=0\quad(r\ne1).
\label{eq:jetF}
\end{equation}
Hence
\begin{equation}
H_m(u):=u-F_m(u)
\label{eq:Hm}
\end{equation}
is flat at $u=0$.
Nevertheless $H_m$ is nonzero in every neighborhood of the origin, and $F_m$
is nowhere real analytic on $\mathbb R$.
\end{theorem}

\begin{proof}
Even derivatives vanish because $F_m$ is odd.  For $r\ge0$,
\[
F_m^{(2r+1)}(0)
=(-1)^r
\sum_{i=0}^{\infty}w^{(m)}_{\infty,i}m^{2ri}.
\]
Proposition~\ref{prop:infinite-moments} gives $F_m'(0)=1$ and zero for every
higher odd derivative.  Thus $H_m$ is flat.

To prove that the flat function is nonzero, first note from
\eqref{eq:limitweights} that for $i\ge1$
\begin{equation}
w^{(m)}_{\infty,i-1}
=-(m^{2i}-1)w^{(m)}_{\infty,i}.
\label{eq:weightratio}
\end{equation}
Using the $C^2$ convergence from Theorem~\ref{thm:Cinf}, termwise
differentiation and an index shift give
\begin{align*}
F_m(mu)
&=m\sum_{i=1}^{\infty}w^{(m)}_{\infty,i-1}
\frac{\sin(m^iu)}{m^i}\\
&=-m\sum_{i=1}^{\infty}(m^{2i}-1)w^{(m)}_{\infty,i}
\frac{\sin(m^iu)}{m^i}\\
&=m\bigl(F_m''(u)+F_m(u)\bigr).
\end{align*}
Thus
\begin{equation}
F_m(mu)=m\bigl(F_m''(u)+F_m(u)\bigr).
\label{eq:funcdiff}
\end{equation}

Suppose $H_m$ vanished on some neighborhood of $0$.  Then $F_m(u)=u$ there,
so $F_m''(u)=0$ there.  Equation \eqref{eq:funcdiff} would then imply
$F_m(v)=v$ on an interval $m$ times larger.  Iterating this argument gives
$F_m(u)=u$ for every real $u$.  But $F_m$ is $2\pi$-periodic, since every term
in \eqref{eq:Fm} is, whereas $u$ is not.  This contradiction proves that
$H_m$ is nonzero in every neighborhood of $0$.

It remains to prove nowhere real analyticity.  Write the nonzero Fourier
coefficient of $F_m$ at frequency $n_i=m^i$ as $\widehat F_m(n_i)$.  From
\eqref{eq:limitweights},
\[
|\widehat F_m(m^i)|
=\frac{|w^{(m)}_{\infty,i}|}{2m^i}
=\frac{m^{-i(i+2)}}{2(q;q)_i(q;q)_\infty}.
\]
Since $(q;q)_i\to(q;q)_\infty>0$,
\begin{equation}
\log |\widehat F_m(m^i)|
=-(i^2+2i)\log m+O(1)
=-\frac{(\log n_i)^2}{\log m}-2\log n_i+O(1).
\label{eq:fourierdecay}
\end{equation}
Thus the Fourier coefficients decay faster than every algebraic power along
the lacunary support, but they do not satisfy an exponential bound
$O(e^{-c|n|})$ for any $c>0$.  A $2\pi$-periodic real-analytic function on
$\mathbb R$ has exponentially decaying Fourier coefficients; see, for example,
Katznelson~\cite[Chapter~I]{Katznelson2004}.  Hence $F_m$ is not real analytic
on all of $\mathbb R$.

Suppose now that $F_m$ were real analytic on some nonempty open interval $I$.
Equation~\eqref{eq:funcdiff} gives, for $v\in mI$,
\[
F_m(v)=m\left(F_m''(v/m)+F_m(v/m)\right),
\]
so $F_m$ is real analytic on the dilated interval $mI$.  Iteration gives
analyticity on $m^kI$ for every $k\ge1$.  For sufficiently large $k$, the
length of $m^kI$ exceeds $2\pi$.  Because $F_m$ is $2\pi$-periodic, analyticity
on such an interval propagates by translation to all of $\mathbb R$,
contradicting the preceding Fourier-coefficient argument.  Therefore $F_m$ is
nowhere real analytic.
\end{proof}

\begin{corollary}[Transported flat limit in the polynomial variable]
\label{cor:zlimit}
Define
\begin{equation}
G_m(z):=F_m(\arcsin z),\qquad -1\le z\le1.
\label{eq:Gm}
\end{equation}
Then $G_m$ is continuous on $[-1,1]$, belongs to $C^\infty((-1,1))$, and
\begin{equation}
P^{(m)}_j(z)\longrightarrow G_m(z)
\end{equation}
uniformly on $[-1,1]$ and in $C^\infty$ on every compact subinterval of
$(-1,1)$.  Moreover,
\begin{equation}
\arcsin z-G_m(z)
\end{equation}
is a nonzero flat function at $z=0$.  In particular, $G_m$ is not real
analytic at the origin despite having the same Taylor series there as
$\arcsin z$.
\end{corollary}

\begin{proof}
Uniform convergence on $[-1,1]$ follows from Theorem~\ref{thm:Cinf} after
putting $u=\arcsin z\in[-\pi/2,\pi/2]$.  On compact subintervals of
$(-1,1)$, composition with the smooth map $\arcsin$ preserves convergence of
all derivatives.  Indeed, with $H_m(u)=u-F_m(u)$ from \eqref{eq:Hm},
\[
\arcsin z-G_m(z)=H_m(\arcsin z).
\]
Flatness follows because $H_m$ is flat at $0$ and $\arcsin z$ is analytic
with $\arcsin z=z+O(z^3)$.  If the flat difference vanished locally in $z$, then
$F_m(u)=u$ locally after setting $z=\sin u$, contradicting
Theorem~\ref{thm:flat}.
\end{proof}

\section{Cubic specialization and interpretation}

For $m=3$ one has $q=1/9$ and
\begin{equation}
w^{(3)}_{\infty,i}
=
\frac{(-1)^i 9^{-i(i+1)/2}}
{(1/9;1/9)_i(1/9;1/9)_\infty}.
\label{eq:cubicweights}
\end{equation}
The finite polynomials are therefore
\begin{equation}
P_j(z)
=
\sum_{i=0}^{j-1}w^{(3)}_{j,i}3^{-i}S_3^{\circ i}(z),
\qquad
S_3(z)=3z-4z^3,
\label{eq:cubiciterate}
\end{equation}
and converge on $[-1,1]$ to
\begin{equation}
G_3(z)
=
\sum_{i=0}^{\infty}w^{(3)}_{\infty,i}
\frac{\sin(3^i\arcsin z)}{3^i}.
\label{eq:cubiclimit}
\end{equation}

For scale, Table~\ref{tab:flatgap} records three evaluations at 50-digit precision (mpmath)
of $H_3(u)=u-F_3(u)$.  These values are only numerical illustrations; the
nonvanishing and flatness statements above are proved independently.
\begin{table}[ht]
\centering
\begin{tabular}{@{}cc@{}}
\toprule
$u$ & $H_3(u)$ \\
\midrule
$0.1$ & $4.7381221300\times10^{-10}$ \\
$0.6$ & $2.2871746237\times10^{-3}$ \\
$1.5$ & $3.1540966034\times10^{-1}$ \\
\bottomrule
\end{tabular}
\caption{Illustrative values of the nonzero flat discrepancy for $m=3$.}
\label{tab:flatgap}
\end{table}
Every finite stage agrees with $\arcsin z$ to an increasing odd order, yet the
infinite-depth limit is not $\arcsin z$.  This resolves the convergence question
left open by the finite-jet theorem: the sequence converges, but the limit
contains a nonperturbative flat discrepancy invisible to every Taylor
coefficient at the fixed point.

This phenomenon clarifies the roles of the different structures.  The
finite-depth coefficient cancellation is classical polynomial extrapolation.
The polynomial composition formula is its exact transport through the
multiple-angle dynamics.  The flat nonanalytic limit appears only after taking
the extrapolation depth to infinity on the outward geometric node set
$1,m^2,m^4,\ldots$.

\section{Scope and concluding remarks}

The finite-depth statements in this paper should be read within classical
Richardson/Romberg theory.  The weights \eqref{eq:finiteweights}, the moment
cancellations \eqref{eq:finite-moments}, and the leading error constant
\eqref{eq:leadingmoment} are standard consequences of polynomial
extrapolation.  Their appearance in \eqref{eq:lift} is explained exactly by
Theorem~\ref{thm:representation}.

The additional structure comes from the multiple-angle transport.  It produces
finite polynomials in $z$ from an extrapolation tableau in $h^2$, gives the
iterate formula \eqref{eq:iterateform}, and makes the same operation appear as
the scale/composition intertwining identity \eqref{eq:intertwining}.  At
infinite depth the geometric Lagrange weights stabilize super-geometrically,
which yields a global $C^\infty$ Hadamard-lacunary trigonometric limit.  The
limiting moment identities force the full Taylor jet of the identity function.
The Fourier-coefficient asymptotic \eqref{eq:fourierdecay}, together with the
functional-differential relation \eqref{eq:funcdiff}, strengthens the local
nonanalyticity statement to nowhere real analyticity.  Thus the limiting
discrepancy is genuinely beyond all algebraic orders at the origin.

No claim is made that Richardson extrapolation, Romberg-type geometric
tableaux, constructive Schr\"oder approximation, Schr\"oder linearization, or
the multiple-angle identity is new.  The paper's claim is the exact structural
transport that turns the classical geometric extrapolation tableau into a
forward finite-polynomial composition hierarchy for $S_m$, together with the
analysis of its infinite-depth limit: a global smooth, nowhere-real-analytic
lacunary trigonometric limit whose transport back to the polynomial variable
has the complete arcsine Taylor jet at the origin but differs from $\arcsin$
by a nonzero flat function.

\begin{thebibliography}{99}

\bibitem{Schroeder1870}
E.~Schr\"oder,
\newblock Ueber iterirte Functionen,
\newblock \emph{Mathematische Annalen} \textbf{3} (1870), no.~2, 296--322.

\bibitem{RichardsonGaunt1927}
L.~F.~Richardson and J.~A.~Gaunt,
\newblock The deferred approach to the limit. Part I. Single lattice. Part II.
Interpenetrating lattices,
\newblock \emph{Philosophical Transactions of the Royal Society of London,
Series A} \textbf{226} (1927), 299--361.
\newblock doi:10.1098/rsta.1927.0008.

\bibitem{Romberg1955}
W.~Romberg,
\newblock Vereinfachte numerische Integration,
\newblock \emph{Norske Videnskabers Selskabs Forhandlinger, Trondheim}
\textbf{28} (1955), 30--36.

\bibitem{Joyce1971}
D.~C.~Joyce,
\newblock Survey of extrapolation processes in numerical analysis,
\newblock \emph{SIAM Review} \textbf{13} (1971), no.~4, 435--490.
\newblock doi:10.1137/1013092.

\bibitem{BrezinskiRedivo1991}
C.~Brezinski and M.~Redivo-Zaglia,
\newblock \emph{Extrapolation Methods: Theory and Practice},
\newblock Studies in Computational Mathematics, vol.~2, North-Holland,
Amsterdam, 1991.

\bibitem{GasperRahman2004}
G.~Gasper and M.~Rahman,
\newblock \emph{Basic Hypergeometric Series}, 2nd ed.,
\newblock Encyclopedia of Mathematics and its Applications, vol.~96,
Cambridge University Press, Cambridge, 2004.

\bibitem{Iserles1993}
A.~Iserles,
\newblock A constructive approach to the Schr\"oder equation,
\newblock \emph{Journal of Computational and Applied Mathematics}
\textbf{46} (1993), 301--314.
\newblock doi:10.1016/0377-0427(93)90304-T.

\bibitem{Sidi2003}
A.~Sidi,
\newblock \emph{Practical Extrapolation Methods: Theory and Applications},
\newblock Cambridge University Press, Cambridge, 2003.
\newblock doi:10.1017/CBO9780511546815.

\bibitem{Zygmund2002}
A.~Zygmund,
\newblock \emph{Trigonometric Series}, vols.~I and II, 3rd ed.,
\newblock Cambridge Mathematical Library, Cambridge University Press,
Cambridge, 2002.

\bibitem{Kahane1964}
J.-P.~Kahane,
\newblock Lacunary Taylor and Fourier series,
\newblock \emph{Bulletin of the American Mathematical Society}
\textbf{70} (1964), 199--213.
\newblock doi:10.1090/S0002-9904-1964-11080-6.

\bibitem{Katznelson2004}
Y.~Katznelson,
\newblock \emph{An Introduction to Harmonic Analysis}, 3rd ed.,
\newblock Cambridge Mathematical Library, Cambridge University Press,
Cambridge, 2004.
\newblock doi:10.1017/CBO9781139165372.

\end{thebibliography}

\end{document}